\chapter{Results and Discussion}
\label{chapter:results_and_discussion}

This chapter presents the empirical evaluation of Quality-Aware Decoding across automatic speech recognition and image captioning tasks. By contrasting conventional Maximum A Posteriori (MAP) decoding with post-hoc reranking methodologies, the analysis details the trade-offs between acoustic/visual grounding, candidate diversity, and computational cost.

\section{ASR Quality Analysis}
\label{sec:asr_quality_analysis}

The evaluation of speech recognition pipelines focuses on the interaction between candidate generation strategies, model size, and decision rules under acoustically degraded conditions. Experiments conducted on the noisy LibriSpeech benchmark illustrate the delicate balance required when selecting candidates from autoregressive sequence generators. When evaluating the lightweight \texttt{whisper-small} architecture under a beam size of 8 using Diverse Beam Search (DBS), standard MAP decoding achieves a baseline Word Error Rate (WER) of 10.12\%[cite: 7]. Applying Minimum Bayes Risk (MBR) decoding over this candidate pool yields a WER of 10.24\%, which perfectly matches the performance of the Two-Stage MBR algorithm[cite: 7]. A parallel evaluation on the high-capacity \texttt{whisper-large-v3} model shows a similar structural pattern: MAP decoding achieves 5.89\% WER at beam size 8 with DBS, while MBR and Two-Stage MBR both settle at 5.93\%[cite: 7]. 

These findings highlight the extreme mode-confidence inherent to heavily trained foundation models. Because architectures like Whisper assign the vast majority of their probability mass to the top-ranked beam, generating additional candidates through stochastic or diverse beam techniques introduces non-trivial risk. On audio signals where the primary acoustic mode identified by MAP decoding is already near-optimal, unweighted consensus operations across alternative candidates risk introducing minor edit penalties rather than correcting transcription errors. However, the exact equivalence between pure MBR and Two-Stage MBR validates the pruning methodology: the pipeline successfully reduces the $O(N^2)$ candidate pool using a fast estimator without incurring any additional degradation in the final WER[cite: 7].

\section{The Referenceless Hallucination Trap in Speech}
\label{sec:hallucination_trap}

The most striking empirical phenomenon observed during ASR decoding occurs when applying fixed referenceless Quality Estimation via the NoRefER metric. While Fixed Reranking maintains stable WER scores under nucleus sampling—yielding 12.49\% for \texttt{whisper-small} at beam size 8—its performance catastrophically degrades when candidate pools are generated using Diverse Beam Search[cite: 7]. Under DBS at beam size 8, the WER surges to an extreme 79.28\% for the small model and 77.29\% for the large model[cite: 7].

This severe degradation provides definitive empirical evidence of the referenceless "hallucination trap" and the pervasive "fluency bias" documented in recent speech evaluation literature. Diverse Beam Search forces the autoregressive decoder to explore structurally distinct paths by applying deterministic penalties across beam groups. In speech models, this aggressive search constraint frequently causes the decoder to decouple from the cross-attention acoustic representations and fall back entirely on its internal language prior, producing fluent, syntactically flawless sentences that bear zero phonetic relationship to the audio input. Because NoRefER is a text-centric metric operating strictly within the textual domain, it evaluates transcripts based on linguistic plausibility and grammatical coherence without any mechanism to verify the underlying acoustic features. Consequently, when presented with a fluent but entirely hallucinated candidate, NoRefER assigns it a deceptively high quality score, forcefully preferring it over acoustically faithful but grammatically complex transcripts. This failure mode demonstrates that referenceless text-only quality estimators are structurally blind to acoustic divergence and cannot be safely deployed as standalone decision rules without acoustic likelihood constraints or strict two-stage candidate pruning.

\section{Image Captioning Quality Analysis}
\label{sec:image_captioning_quality}

The divergence between reference-based consensus metrics and reference-free visual-semantic alignment becomes explicitly clear in the image captioning evaluation. Under a beam size of 5, both Maximum A Posteriori (MAP) decoding and Minimum Bayes Risk (MBR) decoding achieved an identical CIDEr score of 0.8337[cite: 7]. However, applying a Fixed Reranker (F-RR) guided purely by CLIPScore drops the CIDEr performance precipitously to 0.6311[cite: 7]. This phenomenon scales identically at beam size 10, where the baseline MAP score of 0.7961 collapses to 0.5430 under fixed CLIPScore reranking[cite: 7].

This degradation is a mathematically predictable consequence of prioritizing abstract visual grounding over the specific stylistic priors of human language. Vision-Language Models like CLIP operate in a continuous embedding space that frequently behaves as a bag-of-words, rewarding highly specific visual details and keyword chaining even if it shatters the structural flow expected by human annotators. Consequently, generating distinct, image-specific captions naturally incurs a severe penalty under CIDEr, which stringently enforces rigid $n$-gram overlaps. To mitigate this catastrophic drop in syntactical structure, the Two-Stage MBR algorithm successfully bridges the modality gap. By first pruning the candidate pool with CLIPScore to isolate a visually grounded subset, and subsequently applying MBR using CIDEr to identify the most consensus-aligned formulation, the framework recovers significant linguistic structure. At beam size 10, Two-Stage MBR elevated the severely degraded 0.5430 Fixed-RR score back to 0.6893[cite: 7]. This hybrid approach demonstrates that it is possible to constrain reference-free semantic optimization with a structural language prior, retaining visual fidelity while suppressing grammatical degeneration.

\section{Test-Time Compute and Latency Analysis}
\label{sec:test_time_scaling}

The deployment of sequence models introduces a strict tension between inference latency and generation quality. The experimental caching mechanism employed in this pipeline successfully isolates the latency cost of candidate generation from the mathematical overhead of post-hoc reranking. Generating image captioning candidates across a beam size of 5 required 234.96 seconds of compute time, as the autoregressive engine must explicitly generate tokens until termination[cite: 7]. However, executing pure MBR on that generated pool required only 13.44 seconds, and executing Two-Stage MBR required just 30.00 seconds[cite: 7]. 

The latency scaling observed in the speech modality further emphasizes the efficiency of decoupled reranking. Generating candidates for the \texttt{whisper-large-v3} model under Diverse Beam Search at beam size 8 consumed 49,236.12 seconds[cite: 7]. In stark contrast, applying the mathematically complex MBR consensus algorithm over those cached candidates took only 2.82 seconds, while the Two-Stage MBR approach completed the pruning and selection process in 32.30 seconds[cite: 7]. By utilizing fast quality estimators to prune massive candidate pools prior to conducting $O(N^2)$ cross-evaluations, adaptive test-time compute allocation ensures that deep algorithmic search is expended only on probabilistically viable trajectories. This framework drastically minimizes the inference overhead required for complex decision rules, making advanced consensus decoding economically viable for real-world deployment.

\section{Model Scaling and Capability Bridges}
\label{sec:model_scaling}

A primary objective of test-time compute scaling is to ascertain whether smaller, highly optimized models can rival the performance of massive foundation models constrained to standard decoding[cite: 10]. In the noisy ASR experiments, the 244-million parameter \texttt{whisper-small} model, when equipped with Diverse Beam Search and Two-Stage MBR, reduced its baseline error rate from 12.45\% down to 10.12\%[cite: 6]. 

While it does not fully eclipse the 5.87\% WER baseline of the 1.55-billion parameter \texttt{whisper-large-v3} model[cite: 7], the relative efficiency gains achieved by the smaller architecture are profound. The scaling laws observed at inference time suggest that allocating supplementary computational resources toward multi-prompt candidate generation and rigorous consensus verification enables small language models to dynamically stretch their cognitive horizons[cite: 10]. By substituting massive parameter weight scaling with targeted algorithmic deduction, Quality-Aware Decoding establishes a highly viable, economically sustainable paradigm for deploying generative models in precision-critical environments[cite: 10].

\begin{table}[htbp]
\centering
\caption{Automatic Speech Recognition Word Error Rate (WER \%) on the Noisy LibriSpeech benchmark across a candidate pool size of 8. The table contrasts generation algorithms and reranking decision rules, highlighting the severe hallucination vulnerability of referenceless metrics under Diverse Beam Search.}
\label{tab:asr_results_empirical}
\begin{tabular}{lllc}
\toprule
\textbf{Model Architecture} & \textbf{Generation Strategy} & \textbf{Decision Rule} & \textbf{WER (\%)} \\
\midrule
\texttt{whisper-small} (244M) & Nucleus Sampling & MAP & 10.12 \\
\texttt{whisper-small} (244M) & Nucleus Sampling & MBR (WER) & 10.27 \\
\texttt{whisper-small} (244M) & Nucleus Sampling & Fixed-RR (NoRefER) & 12.49 \\
\texttt{whisper-small} (244M) & Nucleus Sampling & Two-Stage MBR & 10.26 \\
\midrule
\texttt{whisper-small} (244M) & Diverse Beam Search & MAP & 10.12 \\
\texttt{whisper-small} (244M) & Diverse Beam Search & MBR (WER) & 10.24 \\
\texttt{whisper-small} (244M) & Diverse Beam Search & Fixed-RR (NoRefER) & \textbf{79.28} \\
\texttt{whisper-small} (244M) & Diverse Beam Search & Two-Stage MBR & 10.24 \\
\midrule
\texttt{whisper-large-v3} (1.55B) & Nucleus Sampling & MAP & 5.89 \\
\texttt{whisper-large-v3} (1.55B) & Nucleus Sampling & MBR (WER) & 5.95 \\
\texttt{whisper-large-v3} (1.55B) & Nucleus Sampling & Fixed-RR (NoRefER) & 6.28 \\
\texttt{whisper-large-v3} (1.55B) & Nucleus Sampling & Two-Stage MBR & 5.95 \\
\midrule
\texttt{whisper-large-v3} (1.55B) & Diverse Beam Search & MAP & 5.89 \\
\texttt{whisper-large-v3} (1.55B) & Diverse Beam Search & MBR (WER) & 5.93 \\
\texttt{whisper-large-v3} (1.55B) & Diverse Beam Search & Fixed-RR (NoRefER) & \textbf{77.29} \\
\texttt{whisper-large-v3} (1.55B) & Diverse Beam Search & Two-Stage MBR & 5.93 \\
\bottomrule
\end{tabular}
\end{table}

\begin{table}[htbp]
\centering
\caption{Image Captioning performance on MS COCO comparing decoding strategies using Salesforce BLIP-Large. The data illustrates the divergence between visual grounding (Fixed-RR) and $n$-gram consensus (MBR), alongside execution latency demonstrating the efficiency of cached reranking operations.}
\label{tab:caption_results_empirical}
\begin{tabular}{clcc}
\toprule
\textbf{Beam Size} & \textbf{Decision Rule} & \textbf{CIDEr} & \textbf{Latency (s)} \\
\midrule
5 & MAP (Baseline) & 0.8337 & 234.96 \\
5 & MBR (CIDEr) & 0.8337 & 13.44 \\
5 & Fixed-RR (CLIPScore) & 0.6311 & 36.02 \\
5 & Two-Stage MBR & 0.6735 & 30.00 \\
\midrule
10 & MAP (Baseline) & 0.7961 & 201.55 \\
10 & MBR (CIDEr) & 0.7961 & 22.80 \\
10 & Fixed-RR (CLIPScore) & 0.5430 & 28.35 \\
10 & Two-Stage MBR & 0.6893 & 27.16 \\
\bottomrule
\end{tabular}
\end{table}
